% Weekly exercise sheet for Lecture 9; also included in the synthesis sheet.
\ifdefined\ExercisesSynthesis
\else
\documentclass[11pt]{article}
\usepackage[utf8]{inputenc}
\usepackage[T1]{fontenc}
\usepackage[english]{babel}
\usepackage{lmodern,microtype}
\usepackage[margin=0.9in]{geometry}
\usepackage{amsmath,amssymb,amsthm,mathtools}
\usepackage{enumitem,hyperref}
\newcommand{\E}{\mathbb E}
\newcommand{\Qbb}{\mathbb Q}
\newcommand{\dd}{\,\mathrm d}
\theoremstyle{definition}
\newtheorem{exercise}{Exercise}[section]
\newlist{exparts}{enumerate}{1}
\setlist[exparts,1]{label=\textbf{(\alph*)},itemsep=0.3em,topsep=0.4em,leftmargin=2em}
\newcommand{\lecturesummary}[1]{\begin{quote}\small\textbf{Focus.} #1\end{quote}}
\title{Stochastic Calculus and Option Pricing\\Lecture 9 Exercises}
\author{}
\date{}
\begin{document}
\maketitle
\vspace{-1.2em}
\setcounter{section}{8}
\fi

\section{Week 9 --- Vanilla Options in Practice}
\lecturesummary{A short theoretical variation on contract units that
remains valid without choosing market numbers.}

\begin{exercise}[A stock split changes units, not economics]
Suppose one old share is replaced by \(a>0\) new shares, so that the quoted
stock price and strike become \(S'=S/a\) and \(K'=K/a\).
\begin{exparts}
  \item Compare the old call payoff \((S_T-K)^+\) with the payoff on \(a\)
  new calls struck at \(K/a\). Explain why the option position must also be
  rescaled by \(a\).
  \item If \(C(S,K)\) is the old call value and \(C'(S',K')\) the value per
  new call, derive the relation between them by no-arbitrage.
  \item Use the chain rule to relate the old and new Delta and Gamma. Which
  sensitivity is invariant to the change of share denomination?
\end{exparts}
\end{exercise}


\ifdefined\ExercisesSynthesis\else
\end{document}
\fi
